\subsection{Open delta}
\label{sec:na-open}

Two single-phase units connected as two sides of a delta supply a
three-phase load, Figure~\ref{fig:na-open}; the third side carries no
winding. With a centre tap on one unit the secondary is a high-leg delta as
in Section~\ref{sec:na-delta}. The connection has no phase shift.

\begin{figure}[!ht]
  \centering
  % generated by north-american-practice/fig/connection_figures.py -- do not edit
\begin{tikzpicture}[x=0.3mm, y=-0.3mm, line cap=round, line join=round]
  \draw[miss] (38.0,130.0) -- (90.0,40.0);
  \node[font=\footnotesize, text=atFgMuted, anchor=center, inner sep=0.5pt] at (51.9,78.0) {open};
  \draw[hv] (90.0,40.0) -- (142.0,130.0);
  \node[font=\footnotesize, text=atInfo, anchor=center, inner sep=0.5pt] at (129.0,77.5) {480};
  \draw[hv] (142.0,130.0) -- (38.0,130.0);
  \node[font=\footnotesize, text=atInfo, anchor=center, inner sep=0.5pt] at (90.0,145.0) {480};
  \fill[atInfo] (38.0,130.0) circle (3.0);
  \node[font=\footnotesize, text=atInfo, anchor=center, inner sep=0.5pt] at (26.8,136.5) {H1};
  \fill[atInfo] (90.0,40.0) circle (3.0);
  \node[font=\footnotesize, text=atInfo, anchor=center, inner sep=0.5pt] at (90.0,27.0) {H2};
  \fill[atInfo] (142.0,130.0) circle (3.0);
  \node[font=\footnotesize, text=atInfo, anchor=center, inner sep=0.5pt] at (153.2,136.5) {H3};
  \draw[miss] (266.0,130.0) -- (318.0,40.0);
  \node[font=\footnotesize, text=atFgMuted, anchor=center, inner sep=0.5pt] at (279.9,78.0) {open};
  \draw[lv] (318.0,40.0) -- (370.0,130.0);
  \node[font=\footnotesize, text=atAccentFg, anchor=center, inner sep=0.5pt] at (357.0,77.5) {240};
  \draw[lv] (370.0,130.0) -- (266.0,130.0);
  \draw[high] (318.0,130.0) -- (318.0,40.0);
  \node[font=\footnotesize, text=atDanger, anchor=east, inner sep=0.5pt] at (313.0,108.0) {208};
  \node[font=\footnotesize, text=atAccentFg, anchor=center, inner sep=0.5pt] at (291.0,122.0) {120};
  \node[font=\footnotesize, text=atAccentFg, anchor=center, inner sep=0.5pt] at (345.0,122.0) {120};
  \fill[atFg] (318.0,130.0) circle (3.0);
  \draw[gnd] (318.0,133.0) -- (318.0,143.0);
  \draw[gnd] (311.0,143.0) -- (325.0,143.0);
  \draw[gnd] (313.5,147.0) -- (322.5,147.0);
  \draw[gnd] (316.0,151.0) -- (320.0,151.0);
  \node[font=\footnotesize, text=atFg, anchor=west, inner sep=0.5pt] at (327.0,146.0) {X0};
  \fill[atAccentFg] (266.0,130.0) circle (3.0);
  \node[font=\footnotesize, text=atAccentFg, anchor=center, inner sep=0.5pt] at (254.8,136.5) {X1};
  \fill[atAccentFg] (318.0,40.0) circle (3.0);
  \node[font=\footnotesize, text=atAccentFg, anchor=center, inner sep=0.5pt] at (318.0,27.0) {X2};
  \fill[atAccentFg] (370.0,130.0) circle (3.0);
  \node[font=\footnotesize, text=atAccentFg, anchor=center, inner sep=0.5pt] at (381.2,136.5) {X3};
  \draw[arr] (172.0,104.0) -- (204.0,104.0);
  \node[font=\footnotesize, text=atFgMuted, anchor=center, inner sep=0.5pt] at (188.0,122.0) {0\textdegree};
  \node[font=\footnotesize, text=atFgMuted, anchor=center, inner sep=0.5pt] at (90.0,196.0) {480\,V, two windings};
  \node[font=\footnotesize, text=atFgMuted, anchor=center, inner sep=0.5pt] at (318.0,196.0) {240/120\,V, two windings};
\end{tikzpicture}

  \caption{Open-delta bank of two single-phase units, secondary with centre
  tap on the winding X3--X1.}
  \label{fig:na-open}
\end{figure}

Each unit carries the full line current. With $S_{1}$ the rating of one unit
and $I_{r} = S_{1} / U_{ll}$ its rated current, the three-phase capacity of
the bank is
\begin{equation}
  S_{v} = \sqrt{3} \cdot U_{ll} \cdot I_{r} = \sqrt{3} \cdot S_{1}
        = 0.866 \cdot (2 \cdot S_{1})
  \label{eq:na-open}
\end{equation}
which is $\sqrt{3}/3 = 57.7$\,\% of the capacity $3 \cdot S_{1}$ of a closed
delta of three units. The two units are loaded asymmetrically, so the
secondary voltages become unbalanced under load. A third unit added later
closes the delta.
\FloatBarrier
